\chapter{Mean-field applications: TDA and RPA}
\label{chap:mfapplications}
%=============================================================================

Chapter~\ref{chap:hartrefock} ended with a matrix.  Expanding the
Hartree-Fock energy to second order in the amplitudes of a Thouless rotation
produced the stability matrix~(\ref{eq:6-stabilitymatrix}), and we remarked in
passing that it is the matrix of the random-phase approximation.  This chapter
makes good on that remark.  We shall see that the same two blocks $A$ and $B$,
arranged Hermitian, decide whether the mean field is a minimum, and arranged
non-Hermitian, give the excitation energies of the system.  Stability and
spectroscopy are two readings of one calculation.

The vehicle is a schematic model with two competing interactions -- a pairing
term that scatters whole pairs between levels, and a particle-hole term that
breaks them.  At zero particle-hole strength it reduces exactly to the pairing
model of Section~\ref{sec:pairingmodel}, so every number can be checked
against Table~\ref{tab:pairing}; switching the particle-hole term on makes the
particle-hole channel non-trivial and lets us watch two mean fields compete.
We solve the model exactly by the methods of Chapter~\ref{chap:fci}, then in
the Tamm-Dancoff approximation, then in the random-phase approximation on the
Hartree-Fock state, and finally in the quasiparticle random-phase
approximation built on a BCS vacuum.  Comparing four approximations against
one exact answer, in a system small enough that nothing is hidden, is the
point of the exercise.

%----------------------------------------------------------------------
\section{The pairing plus particle-hole model}
\label{sec:pairingphmodel}

\index{pairing model!with particle-hole interaction}
We take $L$ doubly degenerate, equally spaced single-particle levels
$p=1,\dots,L$, each carrying a spin label $\sigma=\pm$, occupied by $N$
spin-$\tfrac12$ fermions.  The Hamiltonian is
$\hat{H}=\hat{H}_0+\hat{V}_{\mathrm{pair}}+\hat{V}_{\mathrm{ph}}$ with
\begin{align}
  \hat{H}_0 &= \xi\sum_{p\sigma}(p-1)\,
    \hat{a}^{\dagger}_{p\sigma}\hat{a}_{p\sigma},
  \label{eq:7-H0}\\
  \hat{V}_{\mathrm{pair}} &= -\frac{g}{2}\sum_{pq}
    \hat{a}^{\dagger}_{p+}\hat{a}^{\dagger}_{p-}
    \hat{a}_{q-}\hat{a}_{q+},
  \label{eq:7-Vpair}\\
  \hat{V}_{\mathrm{ph}} &= -\frac{f}{2}\sum_{pqr}
    \left(\hat{a}^{\dagger}_{p+}\hat{a}^{\dagger}_{p-}
      \hat{a}_{q-}\hat{a}_{r+} + \mathrm{h.c.}\right).
  \label{eq:7-Vph}
\end{align}
The one-body term gives level $p$ the energy $\xi(p-1)$; we set $\xi=1$
throughout.  The pairing term, met already in
Section~\ref{sec:pairingmodel}, moves a complete pair -- both spin components
of one level -- from level $q$ to level $p$.  The particle-hole term removes a
spin-down particle from level $q$ and a spin-up particle from level $r$ and
deposits a pair at level $p$; when $q\ne r$ it therefore \emph{breaks} a pair,
which is exactly what the pairing term cannot do.

Both terms conserve $N_+$ and $N_-$ separately, so
$S_z=(N_+-N_-)/2$ is a good quantum number and the low-lying spectrum lives in
the balanced sector $S_z=0$.  For $N=L=4$ the full space has
$\binom{8}{4}=70$ determinants, of which $36$ have $S_z=0$.

Two limits organise everything that follows.

\paragraph{At $f=0$.}
The model is the pairing model of Chapter~\ref{chap:systems}.  The seniority
-- the number of unpaired particles -- is conserved, the ground state has
seniority zero, and the exact energies are those of
Table~\ref{tab:pairing}.  The program \texttt{rpa.py} reproduces them:
\begin{center}
\renewcommand{\arraystretch}{1.2}
\begin{tabular}{rlll}
\toprule
$g$ & $E_{\mathrm{ref}}$ & $E_0$ (this chapter) & Table~\ref{tab:pairing}\\
\midrule
$-1.00$ & $3.000000$ & $2.77987014$ & $2.77987014$\\
$-0.50$ & $2.500000$ & $2.43688426$ & $2.43688426$\\
$0.00$  & $2.000000$ & $2.00000000$ & $2.00000000$\\
$0.50$  & $1.500000$ & $1.41677428$ & $1.41677428$\\
$1.00$  & $1.000000$ & $0.63554847$ & $0.63554847$\\
\bottomrule
\end{tabular}
\end{center}

\paragraph{At $g=0$.}
Only the pair-breaking term survives, and the seniority is maximally violated.
For intermediate $f/g$ the two mechanisms compete, and the interesting
question -- the one the whole chapter is about -- is which mean field to build
the fluctuations on.

%----------------------------------------------------------------------
\section{Exact diagonalisation}
\label{sec:phexact}

The exact solution is Chapter~\ref{chap:fci} applied without modification.
Determinants are bit strings, as in Section~\ref{sec:bitrep}; the Hamiltonian
is built by applying each term of Eqs.~(\ref{eq:7-H0})--(\ref{eq:7-Vph}) to
every basis state and routing the result back to its index with the
accumulated fermion sign; and the lowest eigenvalues follow from the
algorithms of Chapter~\ref{chap:linalg} -- a dense solve when the matrix is
small enough, Lanczos when it is not.

Restricting to $S_z=0$ already cuts the dimension from $70$ to $36$, an
instance of the general remark of Section~\ref{sec:cipractical} that
symmetries are worth exploiting.  For $g=0.5$ and $f=0.05g$ the lowest exact
eigenvalues are
\begin{equation}
  E_0=1.347133,\quad E_1=2.710867,\quad E_2=2.711098,\quad E_3=3.413340,
  \label{eq:7-exactexample}
\end{equation}
with the near-degeneracy of $E_1$ and $E_2$ split only by the small
particle-hole term.  These are the numbers every approximation in this chapter
will be measured against.

%----------------------------------------------------------------------
\section{The Hartree-Fock reference}
\label{sec:phhartreefock}

Both approximations to come are built on a stationary mean field, so we solve
the Hartree-Fock equations of Section~\ref{sec:hfcoefficients} first.  The
antisymmetrised two-body matrix elements are read off the two-particle sector
of $\hat{V}=\hat{H}-\hat{H}_0$,
\begin{equation}
  \bar{v}_{abcd}
   = \left(\hat{a}^{\dagger}_{a}\hat{a}^{\dagger}_{b}\bket{0}\right)^{\dagger}
     \hat{V}\,
     \left(\hat{a}^{\dagger}_{c}\hat{a}^{\dagger}_{d}\bket{0}\right),
  \label{eq:7-vbar}
\end{equation}
and the self-consistent field of Section~\ref{sec:scf} does the rest:

\begin{center}
\renewcommand{\arraystretch}{1.2}
\begin{tabular}{rrrrl}
\toprule
$g$ & $f$ & $E_{\mathrm{HF}}$ & iterations & Hartree-Fock energies\\
\midrule
$0.70$ & $0.00$ & $1.30000$ & $2$ & $-0.350,\ -0.350,\ 0.650,\ 0.650,\dots$\\
$0.70$ & $0.05$ & $1.19709$ & $20$ & $-0.404,\ -0.404,\ 0.600,\ 0.600,\dots$\\
$1.00$ & $0.50$ & $-0.26976$ & $40$ & $-1.505,\ -1.505,\ 0.119,\ 0.119,\dots$\\
\bottomrule
\end{tabular}
\end{center}

The first row is the case of Section~\ref{sec:hfnothing} again: with pure
pairing the filled Fermi sea is already the Hartree-Fock solution, the
iteration converges immediately, and the only effect of the interaction is to
lower each occupied level by $g/2$ while leaving the empty ones alone.  The
particle-hole term does connect the reference to a $1p$-$1h$ state, so the
mean field must rotate the orbitals until Brillouin's theorem,
Eq.~(\ref{eq:6-brillouin}), is restored, and the iteration becomes
non-trivial.

%----------------------------------------------------------------------
\section{The equations of motion}
\label{sec:eom}

\index{equations of motion}
Rather than derive TDA and RPA separately, we obtain both from a single
principle.  Suppose the exact eigenstates satisfy
$\hat{H}\bket{\nu}=E_{\nu}\bket{\nu}$ and define, for each excited state, an
operator $\hat{Q}^{\dagger}_{\nu}$ with
\begin{equation}
  \bket{\nu} = \hat{Q}^{\dagger}_{\nu}\bket{0},
  \qquad
  \hat{Q}_{\nu}\bket{0} = 0,
  \label{eq:7-phonon}
\end{equation}
where $\bket{0}$ is the exact ground state.  Such an operator always exists --
$\hat{Q}^{\dagger}_{\nu}=\ketbra{\nu}{0}$ will do -- and it satisfies
\begin{equation}
  \left[\hat{H},\hat{Q}^{\dagger}_{\nu}\right]\bket{0}
   = \omega_{\nu}\hat{Q}^{\dagger}_{\nu}\bket{0},
  \qquad
  \omega_{\nu} = E_{\nu}-E_0 .
  \label{eq:7-eom}
\end{equation}
Taking the matrix element of Eq.~(\ref{eq:7-eom}) with an arbitrary variation
$\delta\hat{Q}$ and using $\hat{Q}_{\nu}\bket{0}=0$ to symmetrise, one arrives
at the \emph{equation-of-motion} form
\begin{equation}
  \boxed{\;
  \bracket{0}{\left[\delta\hat{Q},\left[\hat{H},
    \hat{Q}^{\dagger}_{\nu}\right]\right]}{0}
   = \omega_{\nu}\,
     \bracket{0}{\left[\delta\hat{Q},
       \hat{Q}^{\dagger}_{\nu}\right]}{0} .
  \;}
  \label{eq:7-eomdouble}
\end{equation}
This is still exact.  It becomes an approximation as soon as we (i) restrict
the operators $\hat{Q}^{\dagger}$ to a manageable family, and (ii) replace the
unknown exact ground state $\bket{0}$ by something we can compute.  Every
method in this chapter is a choice of those two things.

The double commutator on the left has a virtue worth noticing: it is
symmetric in the two operators, so each of them can be evaluated with fewer
contractions than the raw product would need.  This is why the RPA matrices
are traditionally written as double commutators, and it is how the program
computes them -- literally, as expectation values of
$[\hat{O}^{\dagger}_K,[\hat{H},\hat{O}_L]]$ in the chosen reference, with
nothing transcribed by hand.

%----------------------------------------------------------------------
\section{The Tamm-Dancoff approximation}
\label{sec:tda}

\index{Tamm-Dancoff approximation}
The simplest choice is to build the excitation from one particle-hole pair
only, and to take the Hartree-Fock determinant for the ground state:
\begin{equation}
  \hat{Q}^{\dagger}_{\nu} = \sum_{mi}X^{\nu}_{mi}\,
    \hat{a}^{\dagger}_{m}\hat{a}_{i},
  \qquad
  \bket{0}\ \longrightarrow\ \bket{\mathrm{HF}},
  \label{eq:7-tdaoperator}
\end{equation}
with $i,j,\dots$ labelling occupied and $m,n,\dots$ unoccupied Hartree-Fock
orbitals.  Since $\hat{a}^{\dagger}_{m}\hat{a}_{i}$ does annihilate the
determinant when acting to the left, condition~(\ref{eq:7-phonon}) holds
exactly, and the right-hand side of Eq.~(\ref{eq:7-eomdouble}) reduces to
$\delta_{mn}\delta_{ij}$.  What remains is an ordinary Hermitian eigenvalue
problem,
\begin{equation}
  \sum_{nj}A_{mi,nj}X^{\nu}_{nj} = \omega_{\nu}X^{\nu}_{mi},
  \label{eq:7-tdaequation}
\end{equation}
with
\begin{equation}
  A_{mi,nj}
   = \bracket{\mathrm{HF}}{\left[\hat{a}^{\dagger}_{i}\hat{a}_{m},
       \left[\hat{H},\hat{a}^{\dagger}_{n}\hat{a}_{j}\right]\right]}
     {\mathrm{HF}}
   = \left(\varepsilon_m-\varepsilon_i\right)
     \delta_{mn}\delta_{ij} + \bar{v}_{mjin}.
  \label{eq:7-Amatrix}
\end{equation}

There is a plainer way to say what this is.  The matrix $A$ is nothing but the
Hamiltonian restricted to the space of $1p$-$1h$ determinants, measured from
the reference energy,
\begin{equation}
  A_{mi,nj}
   = \bracket{\Phi^m_i}{\hat{H}}{\Phi^n_j}
     - E_{\mathrm{HF}}\,\delta_{mn}\delta_{ij}.
  \label{eq:7-Aisci}
\end{equation}
The program checks this identity directly and finds agreement to
$3\times10^{-15}$.  So the Tamm-Dancoff approximation is exactly the
configuration-interaction truncation of Section~\ref{sec:citruncations}
applied to excited states: keep the reference and all single excitations,
diagonalise, and read off the excitation energies.  It is CIS, no more and no
less, and it inherits both the variational property -- every TDA root is an
upper bound to the corresponding exact excitation energy within that space --
and the defect, that the ground state is left uncorrelated.

%----------------------------------------------------------------------
\section{The random-phase approximation}
\label{sec:rpa}

\index{random-phase approximation}
The Tamm-Dancoff ground state is a single determinant, which we know is wrong.
The random-phase approximation improves it by allowing the phonon to
\emph{de}-excite as well as excite,
\begin{equation}
  \hat{Q}^{\dagger}_{\nu} = \sum_{mi}\left(
    X^{\nu}_{mi}\,\hat{a}^{\dagger}_{m}\hat{a}_{i}
    - Y^{\nu}_{mi}\,\hat{a}^{\dagger}_{i}\hat{a}_{m}\right).
  \label{eq:7-rpaoperator}
\end{equation}
The second term does not annihilate a Slater determinant acting to the right,
so the condition $\hat{Q}_{\nu}\bket{0}=0$ can no longer be satisfied by the
Hartree-Fock state.  It defines instead a new, correlated ground state
$\bket{\mathrm{RPA}}$ -- one that we never construct explicitly, but whose
existence is the whole content of the approximation.

To make progress we need the norm on the right of
Eq.~(\ref{eq:7-eomdouble}), which requires the commutator of two
particle-hole operators,
\begin{equation}
  \left[\hat{a}^{\dagger}_{i}\hat{a}_{m},
    \hat{a}^{\dagger}_{n}\hat{a}_{j}\right]
   = \delta_{mn}\hat{a}^{\dagger}_{i}\hat{a}_{j}
   - \delta_{ij}\hat{a}^{\dagger}_{n}\hat{a}_{m}.
  \label{eq:7-phcommutator}
\end{equation}
This is an operator, not a number.  The \emph{quasiboson approximation}
replaces it by its expectation value in the reference determinant,
\begin{equation}
  \left[\hat{a}^{\dagger}_{i}\hat{a}_{m},
    \hat{a}^{\dagger}_{n}\hat{a}_{j}\right]
  \ \longrightarrow\
  \bracket{\mathrm{HF}}{\left[\hat{a}^{\dagger}_{i}\hat{a}_{m},
    \hat{a}^{\dagger}_{n}\hat{a}_{j}\right]}{\mathrm{HF}}
   = \delta_{mn}\delta_{ij},
  \label{eq:7-quasiboson}
\end{equation}
which is to say that particle-hole pairs are treated as bosons.  It is exact
when the occupations are exactly zero and one and nothing else is going on;
it fails to the extent that the true ground state has particles above the
Fermi level, and it is the origin of every pathology of the method.

With Eq.~(\ref{eq:7-quasiboson}), Eq.~(\ref{eq:7-eomdouble}) becomes the
\emph{RPA eigenvalue problem}
\begin{equation}
  \boxed{\;
  \begin{pmatrix} A & B\\ -B^{*} & -A^{*}\end{pmatrix}
  \begin{pmatrix} X^{\nu}\\ Y^{\nu}\end{pmatrix}
   = \omega_{\nu}
  \begin{pmatrix} X^{\nu}\\ Y^{\nu}\end{pmatrix} ,
  \;}
  \label{eq:7-rpaequation}
\end{equation}
with $A$ as before and
\begin{equation}
  B_{mi,nj}
   = -\bracket{\mathrm{HF}}{\left[\hat{a}^{\dagger}_{i}\hat{a}_{m},
       \left[\hat{H},\hat{a}^{\dagger}_{j}\hat{a}_{n}\right]\right]}
     {\mathrm{HF}}
   = \bar{v}_{mnij}.
  \label{eq:7-Bmatrix}
\end{equation}
The program verifies both closed forms against the double commutators to
better than $10^{-13}$.  Setting $B=0$ discards the $Y$ amplitudes and returns
the Tamm-Dancoff equation, so TDA is RPA with the ground-state correlations
switched off.

Equation~(\ref{eq:7-rpaequation}) is not Hermitian.  Its eigenvalues come in
pairs $\pm\omega_{\nu}$ -- if $(X,Y)$ solves it with $\omega$, then
$(Y^{*},X^{*})$ solves it with $-\omega$ -- and they need not be real.  The
physical solutions are the positive ones, normalised by
\begin{equation}
  \sum_{mi}\left(\vert X^{\nu}_{mi}\vert^2
    - \vert Y^{\nu}_{mi}\vert^2\right) = 1,
  \label{eq:7-rpanorm}
\end{equation}
a metric that is indefinite rather than positive, which is precisely why real
eigenvalues are not guaranteed.

%----------------------------------------------------------------------
\section{RPA and the stability of the mean field}
\label{sec:rpastability}

We can now close the loop with Chapter~\ref{chap:hartrefock}.  The blocks $A$
and $B$ of Eqs.~(\ref{eq:7-Amatrix}) and~(\ref{eq:7-Bmatrix}) are, term for
term, the blocks of the stability matrix~(\ref{eq:6-stabilitymatrix}); only
the arrangement differs,
\begin{equation}
  \hat{M}_{\mathrm{stab}} =
  \begin{pmatrix} A & B\\ B^{*} & A^{*}\end{pmatrix}
  \ \ \text{(Hermitian)},
  \qquad
  \hat{M}_{\mathrm{RPA}} =
  \begin{pmatrix} A & B\\ -B^{*} & -A^{*}\end{pmatrix}
  \ \ \text{(not Hermitian)} .
  \label{eq:7-twoarrangements}
\end{equation}
The relation between them is exact and easy to state.  Writing
$\bm{\eta}=\mathrm{diag}(\bm{1},-\bm{1})$ for the metric of
Eq.~(\ref{eq:7-rpanorm}), we have
$\hat{M}_{\mathrm{RPA}}=\bm{\eta}\,\hat{M}_{\mathrm{stab}}$.  If
$\hat{M}_{\mathrm{stab}}$ is positive definite it has a Hermitian square root,
and
\begin{equation}
  \hat{M}_{\mathrm{stab}}^{1/2}\,\bm{\eta}\,\hat{M}_{\mathrm{stab}}^{1/2}
  \label{eq:7-similarity}
\end{equation}
is Hermitian and similar to $\hat{M}_{\mathrm{RPA}}$, so the RPA frequencies
are real.  Conversely, if $\hat{M}_{\mathrm{stab}}$ has a negative
eigenvalue the argument fails and imaginary roots appear.  Hence

\begin{equation}
  \boxed{\;
  \text{the RPA frequencies are real}
  \iff
  \hat{M}_{\mathrm{stab}}\ge0
  \iff
  \text{the mean field is a local minimum.}
  \;}
  \label{eq:7-stabilityequivalence}
\end{equation}

The Lipkin model of Sections~\ref{sec:lipkin} and~\ref{sec:lipkininstability}
shows this happening.  Feeding the same blocks into both arrangements:

\begin{table}[h]
\centering
\renewcommand{\arraystretch}{1.25}
\setlength{\tabcolsep}{12pt}
\begin{tabular}{rrrl}
\toprule
$\chi$ & $\lambda_{\min}(\hat{M}_{\mathrm{stab}})$
  & imaginary roots & lowest real RPA root\\
\midrule
$0.25$ & $\phantom{-}0.75$ & $0$ & $0.968246$\\
$0.50$ & $\phantom{-}0.50$ & $0$ & $0.866025$\\
$0.90$ & $\phantom{-}0.10$ & $0$ & $0.435890$\\
$1.00$ & $\phantom{-}0.00$ & $0$ & $0.000000$\\
$1.50$ & $-0.50$ & $2$ & ---\\
$2.00$ & $-1.00$ & $2$ & ---\\
\bottomrule
\end{tabular}
\caption{The Lipkin model with $N=4$, $\varepsilon=1$, $W=0$.  The lowest
eigenvalue of the Hermitian stability matrix and the RPA spectrum built from
the same blocks cross zero together at $\chi=1$.  Beyond that the collective
root is imaginary and no excitation energy can be quoted.}
\label{tab:rpastability}
\end{table}

The collective root goes \emph{soft} -- falls to zero -- exactly at $\chi=1$,
where the deformed solution of
Section~\ref{sec:lipkininstability} appears, and becomes imaginary beyond it.
This is the general pattern, and it is one of the most useful things RPA does:
a vanishing frequency announces a mean-field phase transition, and an
imaginary one says the reference is a saddle point.  The excitation spectrum
and the stability analysis are the same calculation read two ways.

For the pairing plus particle-hole model the particle-hole stability matrix
never goes negative -- $\lambda_{\min}$ stays at or above $1$ for every
coupling we consider, including $g=4$, $f=2$.  The particle-hole channel is
simply not where this model becomes collective.  Its instability is in the
pairing channel, and reaching it requires a different mean field, which is the
subject of Sections~\ref{sec:bcs} and~\ref{sec:qrpa}.

%----------------------------------------------------------------------
\section{The RPA ground state and its correlation energy}
\label{sec:rpacorrelation}

Because $\hat{Q}_{\nu}\bket{\mathrm{RPA}}=0$ with $\hat{Q}$ containing the
de-excitation piece, the RPA ground state contains particle-hole pairs.  One
can show, and Exercise~\ref{ex:7-thouless} asks for the argument, that
$\bket{\mathrm{RPA}}$ has the form
\begin{equation}
  \bket{\mathrm{RPA}}
   = \mathcal{N}\exp\left\{\frac{1}{2}\sum_{mnij}
     Z_{mi,nj}\,\hat{a}^{\dagger}_{m}\hat{a}_{i}
       \hat{a}^{\dagger}_{n}\hat{a}_{j}\right\}\bket{\mathrm{HF}},
  \qquad
  \bm{Z} = \bm{Y}^{*}\left(\bm{X}^{*}\right)^{-1},
  \label{eq:7-rpagroundstate}
\end{equation}
an exponential of $2p$-$2h$ operators acting on the Hartree-Fock determinant.
The reader will recognise the structure of Thouless' theorem,
Eq.~(\ref{eq:6-thouless}), with two-body operators in place of one-body ones,
and the structure of the coupled-cluster ansatz,
Eq.~(\ref{eq:5-ccsd}), with the doubles amplitude fixed by the RPA amplitudes
rather than solved for independently.  RPA is a doubles theory in disguise.

Its energy relative to Hartree-Fock is
\begin{equation}
  E^{\mathrm{RPA}}_{\mathrm{corr}}
   = \frac{1}{2}\left(\sum_{\nu}\omega_{\nu} - \mathrm{Tr}\,A\right),
  \label{eq:7-rpacorrelation}
\end{equation}
a remarkably compact result: the correlation energy is half the difference
between the sum of the RPA frequencies and the sum of the TDA ones.  Since
each RPA root lies below the corresponding TDA root, the correction is always
negative -- RPA always lowers the energy.  For $g=0.7$, $f=0.05$ the program
gives $\mathrm{Tr}\,A=38.423313$ and $\sum_{\nu}\omega_{\nu}=37.824314$, hence
$E_{\mathrm{corr}}=-0.299500$.

%----------------------------------------------------------------------
\section{Results in the particle-hole channel}
\label{sec:phresults}

We can now compare.  The table gives, for four combinations of the two
couplings, the Hartree-Fock energy, the RPA-corrected ground-state energy, the
exact ground-state energy, and the lowest excitation energy in the three
treatments.

\begin{table}[h]
\centering
\renewcommand{\arraystretch}{1.25}
\setlength{\tabcolsep}{7pt}
\begin{tabular}{rrrrrrrr}
\toprule
$g$ & $f$ & $E_{\mathrm{HF}}$
  & $E_{\mathrm{HF}}+E_{\mathrm{corr}}$ & $E_{\mathrm{exact}}$
  & $\omega_1^{\mathrm{TDA}}$ & $\omega_1^{\mathrm{RPA}}$
  & $\omega_1^{\mathrm{exact}}$\\
\midrule
$0.50$ & $0.025$ & $1.44927$ & $1.29898$ & $1.34713$
  & $1.2748$ & $1.2448$ & $1.3637$\\
$0.70$ & $0.050$ & $1.19709$ & $0.89759$ & $0.96976$
  & $1.3992$ & $1.3412$ & $1.5972$\\
$1.00$ & $0.050$ & $0.89709$ & $0.36862$ & $0.45058$
  & $1.5492$ & $1.4488$ & $1.9425$\\
$1.00$ & $0.500$ & $-0.26976$ & $-1.73461$ & $-1.69174$
  & $1.8111$ & $1.6445$ & $2.8596$\\
\bottomrule
\end{tabular}
\caption{The pairing plus particle-hole model with $N=L=4$, $\xi=1$.
Hartree-Fock, RPA and exact ground-state energies, and the lowest excitation
energy in the Tamm-Dancoff and random-phase approximations compared with the
exact one.}
\label{tab:phresults}
\end{table}

Three things are worth reading off.

First, RPA always lowers the energy below Hartree-Fock, by construction, and
here it \emph{overshoots} the exact answer in every row.  This is the standard
behaviour: the quasiboson approximation, by treating particle-hole pairs as
bosons that can be created without limit, over-counts the ground-state
correlations.  The overshoot grows with the coupling, from $0.048$ at $g=0.5$
to $0.043$ at $g=1$, $f=0.5$ -- modest here, but in a strongly collective
system RPA can overbind badly and eventually collapse altogether.

Second, $\omega_1^{\mathrm{RPA}}<\omega_1^{\mathrm{TDA}}$ in every row.  The
$B$ block always pushes the collective root down, for the same reason it
lowers the ground state: it adds correlations to the ground state, and the
excitation energy is measured from a lower floor.

Third, both approximations put the lowest excitation well below the exact
one, and the gap widens as the pairing strength grows -- at $g=1$, $f=0.5$
the exact first excitation is $2.86$ against an RPA value of $1.64$.  This is
not a failure of arithmetic but of the choice of channel.  The particle-hole
phonon knows nothing about pairing correlations, and in this model the
low-lying spectrum is shaped by them.

\paragraph{A closed form.}
For pure pairing the particle-hole problem collapses.  All sixteen
particle-hole pairs are degenerate, the matrices reduce to multiples of simple
structures, and one finds
\begin{equation}
  \omega^{\mathrm{TDA}} = \xi + \frac{g}{2},
  \qquad
  \omega^{\mathrm{RPA}} = \sqrt{\xi^2+g\xi}
   = \sqrt{\left(\xi+\tfrac{g}{2}\right)^2-\left(\tfrac{g}{2}\right)^2},
  \label{eq:7-closedform}
\end{equation}
the second being the classic $\sqrt{A^2-B^2}$ of a two-level RPA problem with
$A=\xi+g/2$ and $B=g/2$.  The program confirms both to eight digits:

\begin{center}
\renewcommand{\arraystretch}{1.2}
\begin{tabular}{rllll}
\toprule
$g$ & TDA & $\xi+g/2$ & RPA & $\sqrt{1+g}$\\
\midrule
$0.5$ & $1.25000000$ & $1.25000000$ & $1.22474487$ & $1.22474487$\\
$1.0$ & $1.50000000$ & $1.50000000$ & $1.41421356$ & $1.41421356$\\
$2.0$ & $2.00000000$ & $2.00000000$ & $1.73205081$ & $1.73205081$\\
$3.0$ & $2.50000000$ & $2.50000000$ & $2.00000000$ & $2.00000000$\\
$5.0$ & $3.50000000$ & $3.50000000$ & $2.44948974$ & $2.44948974$\\
\bottomrule
\end{tabular}
\end{center}

Equation~(\ref{eq:7-closedform}) is worth staring at.  The RPA root is real
for all $g>-\xi$, so the particle-hole channel of the pure pairing model never
goes unstable, however strong the pairing.  The instability is elsewhere, and
finding it requires a mean field that can break a different symmetry.

%----------------------------------------------------------------------
\section{BCS: the pairing mean field}
\label{sec:bcs}

\index{BCS theory}\index{pairing gap}
The pairing interaction is
$-G\sum_{pq}\hat{P}^{\dagger}_{p}\hat{P}_{q}$ with
$\hat{P}^{\dagger}_{p}=\hat{a}^{\dagger}_{p+}\hat{a}^{\dagger}_{p-}$ and
$G=g/2$.  Its natural mean field is not of Hartree-Fock type at all, because
the object that acquires an expectation value is
$\langle\hat{P}^{\dagger}_{p}\rangle$, which does not conserve particle
number.  The trial state is the BCS wave function
\begin{equation}
  \bket{\mathrm{BCS}} = \prod_{p}\left(u_p + v_p\,
    \hat{a}^{\dagger}_{p+}\hat{a}^{\dagger}_{p-}\right)\bket{0},
  \qquad u_p^2+v_p^2 = 1,
  \label{eq:7-bcs}
\end{equation}
a product over levels in which each level is empty with amplitude $u_p$ and
occupied by a pair with amplitude $v_p$.  It is a superposition of states with
different particle numbers, so the number is fixed only on average, by
minimising $\bracket{\mathrm{BCS}}{\hat{H}-\lambda\hat{N}}{\mathrm{BCS}}$ with
$\lambda$ adjusted until $\langle\hat{N}\rangle=N$.

Because the trial state is a product over levels, every expectation value is
available in closed form.  Writing $u_p=\cos\theta_p$, $v_p=\sin\theta_p$,
\begin{equation}
  \langle\hat{N}\rangle = 2\sum_p v_p^2,
  \qquad
  \langle\hat{H}_0\rangle = 2\sum_p \varepsilon_p v_p^2,
  \label{eq:7-bcsN}
\end{equation}
and, using $\langle\hat{P}^{\dagger}_p\hat{P}_q\rangle = u_pv_pu_qv_q$ for
$p\ne q$ and $v_p^2$ for $p=q$,
\begin{equation}
  \langle\hat{H}\rangle
   = 2\sum_p\varepsilon_p v_p^2
   - S\left[\sum_p v_p^2
     + \sum_{p\ne q}u_pv_pu_qv_q\right],
  \qquad
  S = \frac{g}{2} + f .
  \label{eq:7-bcsenergy}
\end{equation}
The program checks Eq.~(\ref{eq:7-bcsenergy}) against the expectation value
computed directly in the $2^{2L}$-dimensional Fock space and finds agreement
to $5\times10^{-16}$.

\paragraph{The particle-hole term as extra pairing.}
Equation~(\ref{eq:7-bcsenergy}) contains a small surprise.  Only the $q=r$
part of $\hat{V}_{\mathrm{ph}}$ survives in the BCS state -- for $q\ne r$ the
operator $\hat{a}_{q-}\hat{a}_{r+}$ leaves two broken pairs, and the BCS state
has none -- and $\hat{a}_{q-}\hat{a}_{q+}=\hat{P}_q$ exactly.  Seen from the
BCS vacuum, therefore, the particle-hole term \emph{is} a pairing term, and
the whole problem depends on $g$ and $f$ only through the combination
$g+2f$.  The program confirms it:

\begin{center}
\renewcommand{\arraystretch}{1.2}
\begin{tabular}{rrrll}
\toprule
$g$ & $f$ & $g+2f$ & $E_{\mathrm{BCS}}$ & gap $\Delta$\\
\midrule
$1.4$ & $0.0$ & $1.4$ & $0.515684$ & $0.789218$\\
$1.2$ & $0.1$ & $1.4$ & $0.515684$ & $0.789218$\\
$0.8$ & $0.3$ & $1.4$ & $0.515684$ & $0.789218$\\
$0.0$ & $0.7$ & $1.4$ & $0.515684$ & $0.789218$\\
\bottomrule
\end{tabular}
\end{center}

This is a clean illustration of a general point: a mean field sees only that
part of the interaction which can be written as a one-body field in its own
channel.  Hartree-Fock and BCS decompose the same $\hat{V}$ differently, and
each is blind to what the other sees.

\paragraph{The pairing transition.}
With $\Delta=S\sum_p u_pv_p$ the pairing gap, the solution of the BCS
equations for $f=0$ is

\begin{table}[h]
\centering
\renewcommand{\arraystretch}{1.25}
\setlength{\tabcolsep}{10pt}
\begin{tabular}{rrrl}
\toprule
$g$ & $\Delta$ & $\lambda$ & $v_p^2$\\
\midrule
$0.40$ & $0.00000$ & $1.2500$ & $1.000,\ 1.000,\ 0.000,\ 0.000$\\
$0.80$ & $0.00000$ & $1.2500$ & $1.000,\ 1.000,\ 0.000,\ 0.000$\\
$1.00$ & $0.00000$ & $1.2500$ & $1.000,\ 1.000,\ 0.000,\ 0.000$\\
$1.20$ & $0.46716$ & $1.2000$ & $0.984,\ 0.925,\ 0.075,\ 0.016$\\
$1.60$ & $1.05608$ & $1.1000$ & $0.934,\ 0.783,\ 0.217,\ 0.066$\\
$2.00$ & $1.54088$ & $1.0000$ & $0.887,\ 0.709,\ 0.291,\ 0.113$\\
\bottomrule
\end{tabular}
\caption{The BCS solution of the pairing model with $L=4$ levels and $N=4$
particles.  Below $g_c=1.07037$ the only solution is the normal Fermi sea;
above it a gap opens and the occupations become fractional.}
\label{tab:bcsgap}
\end{table}

Below a critical coupling the only solution is the normal Fermi sea, with
$\Delta=0$ and sharp occupations; above it a gap opens and the occupation
numbers spread.  The condition for a non-trivial solution, obtained by
linearising the gap equation about $\Delta=0$, is
\begin{equation}
  1 = 2S\left[\frac{1}{3\xi+S}+\frac{1}{\xi+S}\right],
  \label{eq:7-gc}
\end{equation}
whose root is $S_c=0.53519$, so that $g_c=1.07037$ at $f=0$ -- exactly the
value the program locates by bisection.

Two comments.  The transition is \emph{sharp} even though the system is
finite, but that sharpness is an artefact of the mean field: the exact
ground-state energy of Table~\ref{tab:pairing} is perfectly smooth through
$g_c$, and nothing whatever happens there.  Mean-field phase transitions in
finite systems are approximations to smooth crossovers.  And above $g_c$ the
reference has stopped being an eigenstate of $\hat{N}$, which is a broken
symmetry, with all that this implies for the fluctuations built on top of it.

%----------------------------------------------------------------------
\section{Quasiparticle RPA}
\label{sec:qrpa}

\index{quasiparticle random-phase approximation}\index{QRPA}
The BCS state is a vacuum, not for particles, but for the Bogoliubov
quasiparticles
\begin{equation}
  \hat{\alpha}^{\dagger}_{p+} = u_p\hat{a}^{\dagger}_{p+}
    - v_p\hat{a}_{p-},
  \qquad
  \hat{\alpha}^{\dagger}_{p-} = u_p\hat{a}^{\dagger}_{p-}
    + v_p\hat{a}_{p+},
  \label{eq:7-bogoliubov}
\end{equation}
which satisfy $\hat{\alpha}_{p\sigma}\bket{\mathrm{BCS}}=0$.  The program
verifies this to $3\times10^{-17}$.  These are the natural elementary
excitations of the paired system, and the quasiparticle random-phase
approximation builds collective modes out of \emph{pairs} of them,
\begin{equation}
  \hat{Q}^{\dagger}_{\nu} = \sum_{a<b}\left(
    X^{\nu}_{ab}\,\hat{\alpha}^{\dagger}_{a}\hat{\alpha}^{\dagger}_{b}
    - Y^{\nu}_{ab}\,\hat{\alpha}_{b}\hat{\alpha}_{a}\right).
  \label{eq:7-qrpaoperator}
\end{equation}
The structure of Eq.~(\ref{eq:7-eomdouble}) is unchanged; only the reference
and the elementary operators are different, so the same routine produces
\begin{equation}
  A_{ab,cd} = \bracket{\mathrm{BCS}}{\left[
    \hat{\alpha}_{b}\hat{\alpha}_{a},
    \left[\hat{H}',
      \hat{\alpha}^{\dagger}_{c}\hat{\alpha}^{\dagger}_{d}\right]\right]}
    {\mathrm{BCS}},
  \qquad
  B_{ab,cd} = -\bracket{\mathrm{BCS}}{\left[
    \hat{\alpha}_{b}\hat{\alpha}_{a},
    \left[\hat{H}',
      \hat{\alpha}_{d}\hat{\alpha}_{c}\right]\right]}{\mathrm{BCS}},
  \label{eq:7-qrpaAB}
\end{equation}
and the same eigenvalue problem~(\ref{eq:7-rpaequation}).  For $L=4$ there are
$2L=8$ quasiparticle states and hence $2L(2L-1)/2=28$ two-quasiparticle
operators.

Two features require attention, and they are the physics of this section.

\subsection{The spurious mode}
\label{sec:spurious}

\index{spurious state}\index{Goldstone mode}
The BCS state breaks the $U(1)$ symmetry generated by $\hat{N}$.  Whenever a
mean field breaks a continuous symmetry of the Hamiltonian, the generator of
that symmetry produces an exact zero mode of the RPA equations -- a Goldstone
mode.  The argument is short: since $[\hat{H},\hat{N}]=0$, the amplitudes
\begin{equation}
  X^{\mathrm{sp}}_{K} = \bracket{\mathrm{BCS}}{
    \left[\hat{O}_K,\hat{N}\right]}{\mathrm{BCS}},
  \qquad
  Y^{\mathrm{sp}}_{K} = \bracket{\mathrm{BCS}}{
    \left[\hat{O}^{\dagger}_K,\hat{N}\right]}{\mathrm{BCS}}
  \label{eq:7-spurious}
\end{equation}
satisfy $\hat{M}_{\mathrm{RPA}}\,(X^{\mathrm{sp}},Y^{\mathrm{sp}})^{T}=0$.
Physically the mode is not an excitation at all: it is the rotation of the
BCS phase, which costs nothing and merely moves the system to a state with a
different average particle number.

The standard device is to build the QRPA with the \emph{cranked} Hamiltonian
\begin{equation}
  \hat{H}' = \hat{H} - \lambda\hat{N},
  \label{eq:7-cranked}
\end{equation}
using the same $\lambda$ that the BCS equations extremise, which places the
spurious root at zero where it can be discarded.  In finite arithmetic it does
not land exactly on zero, and because the RPA metric is singular there it may
emerge as a small \emph{imaginary} pair rather than a small real one.  This is
worth emphasising, since it is easily mistaken for an instability:

\begin{table}[h]
\centering
\renewcommand{\arraystretch}{1.25}
\setlength{\tabcolsep}{11pt}
\begin{tabular}{rrrrr}
\toprule
$g$ & $\Delta$ & $\vert\omega_{\mathrm{sp}}\vert$
  & overlap with $\hat{N}$ & imaginary roots\\
\midrule
$1.20$ & $0.46716$ & $1.14\times10^{-4}$ & $1.000000$ & $0$\\
$1.40$ & $0.78922$ & $1.83\times10^{-4}$ & $1.000000$ & $0$\\
$1.60$ & $1.05608$ & $2.59\times10^{-4}$ & $1.000000$ & $0$\\
$1.80$ & $1.30338$ & $1.64\times10^{-4}$ & $1.000000$ & $0$\\
$2.00$ & $1.54088$ & $2.77\times10^{-4}$ & $1.000000$ & $0$\\
\bottomrule
\end{tabular}
\caption{The spurious mode of the QRPA for the pairing model, $f=0$.  The
root nearest zero has unit overlap with the amplitudes~(\ref{eq:7-spurious})
built from the number operator, identifying it beyond doubt.  Once it is
removed, no imaginary eigenvalues remain at any coupling.}
\label{tab:spurious}
\end{table}

The overlap with the number-operator amplitudes is one to six decimal places,
which identifies the near-zero root beyond any doubt.  Once it is removed, the
QRPA matrix has no imaginary eigenvalues anywhere in the superfluid regime.
The lesson generalises: before declaring a mean field unstable on the strength
of an imaginary RPA root, check whether the root is a spurious mode in
disguise.  The test is cheap -- compute the amplitudes~(\ref{eq:7-spurious})
for each broken symmetry and take overlaps.

\subsection{The mode content}
\label{sec:modecontent}

The second feature is that quasiparticles mix particle number, so a
two-quasiparticle phonon need not conserve it.  Building the phonon operator
explicitly, applying it to the BCS vacuum and taking the expectation value of
$\hat{N}$ separates the modes cleanly:

\begin{table}[h]
\centering
\renewcommand{\arraystretch}{1.25}
\setlength{\tabcolsep}{11pt}
\begin{tabular}{rrl}
\toprule
$\omega_{\nu}$ & $\langle\Delta\hat{N}\rangle$ & character\\
\midrule
\multicolumn{3}{l}{\emph{$g=2.0$, $f=0$; exact excitations
  $3.3274$, $3.3274$, $3.4897$, $4.2536$}}\\
$2.6015$ & $0.0000$ & excitation of the $N$-particle system\\
$2.9212$ & $0.0000$ & excitation of the $N$-particle system (fourfold)\\
$3.7879$ & $-0.3566$ & pairing vibration\\
$3.7879$ & $+0.3566$ & pairing vibration\\
\bottomrule
\end{tabular}
\caption{QRPA phonons at $g=2$, classified by the average particle number
they carry.  Modes with $\langle\Delta\hat{N}\rangle=0$ are excitations
within the $N$-particle system; the others connect towards $N\pm2$ and are
the pairing vibrations.  The shift is not quantised because the BCS vacuum is
not an eigenstate of $\hat{N}$.}
\label{tab:modecontent}
\end{table}

Only the $\langle\Delta\hat{N}\rangle=0$ modes are directly comparable with
the exact spectrum of the $N$-particle system; the others describe transitions
to the neighbouring even systems and would be compared with $E_0(N\pm2)$.
Notice that even the comparable ones lie well below the exact excitations --
$2.60$ and $2.92$ against $3.33$ and $3.49$ -- the same downward bias RPA
showed in the particle-hole channel.

%----------------------------------------------------------------------
\section{Everything against everything}
\label{sec:allresults}

We can now put the four treatments side by side.  For pure pairing, where the
exact answers are those of Chapter~\ref{chap:systems}:

\begin{table}[h]
\centering
\renewcommand{\arraystretch}{1.25}
\setlength{\tabcolsep}{9pt}
\begin{tabular}{rrrrrr}
\toprule
$g$ & $E_{\mathrm{exact}}$ & $E_{\mathrm{HF}}$
  & $E_{\mathrm{HF}}+E^{\mathrm{RPA}}_{\mathrm{corr}}$
  & $E_{\mathrm{BCS}}$
  & $E_{\mathrm{BCS}}+E^{\mathrm{QRPA}}_{\mathrm{corr}}$\\
\midrule
$0.40$ & $\phantom{-}1.548001$ & $1.60000$ & $\phantom{-}1.51748$
  & $\phantom{-}1.60000$ & $\phantom{-}1.46366$\\
$0.60$ & $\phantom{-}1.277449$ & $1.40000$ & $\phantom{-}1.22390$
  & $\phantom{-}1.40000$ & $\phantom{-}1.08964$\\
$0.80$ & $\phantom{-}0.973552$ & $1.20000$ & $\phantom{-}0.90179$
  & $\phantom{-}1.20000$ & $\phantom{-}0.62540$\\
$1.00$ & $\phantom{-}0.635548$ & $1.00000$ & $\phantom{-}0.55459$
  & $\phantom{-}1.00000$ & $-0.01157$\\
$1.20$ & $\phantom{-}0.264504$ & $0.80000$ & $\phantom{-}0.18506$
  & $\phantom{-}0.78522$ & $-0.40860$\\
$1.40$ & $-0.137035$ & $0.60000$ & $-0.20453$
  & $\phantom{-}0.51568$ & $-0.67091$\\
$1.60$ & $-0.565660$ & $0.40000$ & $-0.61230$
  & $\phantom{-}0.20452$ & $-1.04351$\\
$2.00$ & $-1.489652$ & $0.00000$ & $-1.47622$
  & $-0.50353$ & $-1.94223$\\
\bottomrule
\end{tabular}
\caption{Ground-state energies for the pure pairing model, $N=L=4$, $\xi=1$.
The Hartree-Fock energy is the unperturbed reference, since the pairing
interaction cannot couple it to a $1p$-$1h$ state.  The BCS energy departs
from it only above $g_c=1.07$.}
\label{tab:allresults}
\end{table}

The table repays reading across.

\emph{Hartree-Fock} does nothing at all, exactly as
Section~\ref{sec:hfnothing} predicted: $E_{\mathrm{HF}}=2-g$ is the
unperturbed reference energy for every $g$.  The whole correlation energy has
to come from the fluctuations.

\emph{Particle-hole RPA} supplies almost all of it, and tracks the exact curve
remarkably closely -- $-0.61230$ against $-0.565660$ at $g=1.6$ -- crossing it
between $g=1.6$ and $g=2$.  It is slightly under-bound at weak coupling and
slightly over-bound at strong coupling.  That a particle-hole theory should do
this well on a pure pairing interaction is a reminder that the labels
``particle-hole'' and ``pairing'' refer to how a calculation is organised, not
to a partition of the physics.

\emph{BCS} improves on Hartree-Fock only above $g_c$, and then only modestly.
A mean field that merely smears the occupation numbers cannot reproduce a
genuinely correlated ground state; at $g=2$ it recovers $0.50$ of the required
$1.49$.

\emph{QRPA} on top of BCS overshoots badly, and
Section~\ref{sec:modecontent} says why: the two-quasiparticle space contains
pairing vibrations reaching towards $N\pm2$, and the standard
formula~(\ref{eq:7-rpacorrelation}) counts every root.  Restoring particle
number by projection is necessary before BCS plus QRPA becomes quantitative
for a system this small.  In a heavy nucleus, with many levels and a small
relative number fluctuation, the same machinery is far better behaved -- which
is why it remains a standard tool there.

\begin{notebox}
\textbf{What to take away.}  Both approximations of this chapter are the
\emph{same} approximation applied to different reference states and different
elementary operators: linearise the equations of motion, replace the
commutator of two elementary excitations by its expectation value in the
reference, and diagonalise.  Everything that follows -- the $\pm\omega$
pairing of the roots, the indefinite metric, the possibility of imaginary
frequencies, the correlation-energy formula, the appearance of Goldstone modes
whenever the reference breaks a symmetry -- follows from that one step, and is
independent of which channel one is working in.  The quality of the answer,
however, is not: it depends entirely on whether the reference has already
captured the physics, and the two halves of Table~\ref{tab:allresults} show
the same interaction treated well and badly by that criterion.
\end{notebox}

%----------------------------------------------------------------------
\section{Summary}
\label{sec:ch7summary}

The Tamm-Dancoff approximation diagonalises the Hamiltonian in the
$1p$-$1h$ space; it is the excited-state version of the CIS truncation of
Chapter~\ref{chap:fci}, and its matrix $A$ is the Hamiltonian restricted to
that space, measured from the reference.  The random-phase approximation adds
the de-excitation amplitudes $Y$, which correlates the ground state, lowers
every root, and turns the Hermitian eigenvalue problem into the non-Hermitian
one of Eq.~(\ref{eq:7-rpaequation}).  The blocks $A$ and $B$ that appear there
are precisely the blocks of Chapter~\ref{chap:hartrefock}'s stability matrix,
and the RPA frequencies are real exactly when the mean field is a local
minimum: excitation spectrum and stability analysis are one calculation.

Applied to a model with competing pairing and particle-hole interactions, the
approximations behave as their construction demands.  Particle-hole RPA is
stable at every coupling and does surprisingly well on the ground-state
energy, while systematically underestimating the excitation energies.  BCS
finds the pairing instability that the particle-hole channel cannot see, at
the price of breaking particle-number conservation, and QRPA built on it
carries a spurious Goldstone mode which must be identified and removed --
by its overlap with the number-operator amplitudes, not by inspecting the
spectrum alone.

The program \texttt{rpa.py} performs every calculation quoted here, and the
notebook \texttt{tdarpa.ipynb} runs it cell by cell.  All the matrices are
built as literal double commutators, so the same routine serves the
particle-hole and the quasiparticle case, and the textbook formulas
(\ref{eq:7-Amatrix}) and (\ref{eq:7-Bmatrix}) are checked against them rather
than assumed.

%----------------------------------------------------------------------
\section{Exercises}
\label{sec:ch7exercises}
%=============================================================================

\subsection*{Exercise 1: the model and its limits}
\label{ex:7-model}
\addcontentsline{toc}{subsection}{Exercise 1: the model and its limits}

\begin{enumerate}
  \item[a)] Show that $\hat{V}_{\mathrm{pair}}$ and $\hat{V}_{\mathrm{ph}}$
  both conserve $N_+$ and $N_-$ separately, and hence that $S_z$ is a good
  quantum number.  What is the dimension of the $S_z=0$ sector for $N=L=4$?
  \item[b)] Show that at $f=0$ the seniority is conserved, and that at
  $f\ne0$ it is not.  Which matrix elements does the particle-hole term
  supply that the pairing term does not?
  \item[c)] Verify with \texttt{rpa.py} that the $f=0$ spectrum reproduces
  Table~\ref{tab:pairing}, and that the ground state at $f=0$ lies entirely
  in the seniority-zero subspace.
  \item[d)] For $g=0.5$, $f=0.05g$ the two lowest excited states are split by
  only $2\times10^{-4}$.  Explain the near-degeneracy.
\end{enumerate}

\paragraph{Answer to a).}
Every term of Eq.~(\ref{eq:7-Vpair}) contains one $\hat{a}^{\dagger}_{p+}$ and
one $\hat{a}_{q+}$, so $N_+$ is unchanged, and likewise for $N_-$.  In
Eq.~(\ref{eq:7-Vph}) the operator
$\hat{a}^{\dagger}_{p+}\hat{a}^{\dagger}_{p-}\hat{a}_{q-}\hat{a}_{r+}$ again
creates and destroys one particle of each spin.  Hence
$[\hat{H},\hat{N}_{\pm}]=0$ and $S_z=(N_+-N_-)/2$ is conserved.  For $N=L=4$
the balanced sector requires two spin-up and two spin-down particles among
four levels each, so its dimension is $\binom{4}{2}^2=36$, against
$\binom{8}{4}=70$ for the full space.

\paragraph{Answer to b).}
The pairing term moves both members of a pair together, so the number of
unpaired particles -- the seniority -- is unchanged, as
Exercise~\ref{ex:5-blocks}c) showed.  The particle-hole term with $q\ne r$
removes a spin-down from level $q$ and a spin-up from a \emph{different}
level $r$, leaving unpaired particles behind at both, and deposits a pair at
$p$.  It therefore changes the seniority by two.  These are exactly the matrix
elements that connect the seniority-zero subspace to the rest of the $S_z=0$
space, and they are why the full $36$-dimensional problem must be solved once
$f\ne0$.

\paragraph{Answer to c).}
Running \texttt{demo\_model()} prints the $f=0$ energies alongside
Table~\ref{tab:pairing}; they agree to all eight digits shown.  Projecting the
ground-state eigenvector onto the determinants with broken pairs gives zero to
machine precision.

\paragraph{Answer to d).}
At $f=0$ the first excited state of the seniority-zero problem is doubly
degenerate by the symmetry of the level scheme.  The particle-hole term lifts
the degeneracy, but only at order $f$, and $f=0.025$ here; the splitting
$2.711098-2.710867=2.3\times10^{-4}$ is of order $f^2/\xi$, consistent with
the degeneracy being lifted at second order.

\subsection*{Exercise 2: from the equations of motion to TDA and RPA}
\label{ex:7-eom}
\addcontentsline{toc}{subsection}{Exercise 2: from the equations of motion to TDA and RPA}

\begin{enumerate}
  \item[a)] Derive Eq.~(\ref{eq:7-eomdouble}) from Eq.~(\ref{eq:7-eom}),
  stating where $\hat{Q}_{\nu}\bket{0}=0$ is used.
  \item[b)] Show that with the operator~(\ref{eq:7-tdaoperator}) and a Slater
  determinant reference, Eq.~(\ref{eq:7-eomdouble}) reduces to the Hermitian
  problem~(\ref{eq:7-tdaequation}), and evaluate
  $A_{mi,nj}$ using Wick's theorem.
  \item[c)] Show that the Tamm-Dancoff matrix is the Hamiltonian restricted
  to the $1p$-$1h$ space, Eq.~(\ref{eq:7-Aisci}), and hence that TDA is CIS.
  \item[d)] With the operator~(\ref{eq:7-rpaoperator}), derive
  Eq.~(\ref{eq:7-rpaequation}) and identify precisely where the quasiboson
  approximation enters.
  \item[e)] Show that if $(X,Y)$ solves Eq.~(\ref{eq:7-rpaequation}) with
  frequency $\omega$, then $(Y^{*},X^{*})$ solves it with $-\omega$.
\end{enumerate}

\paragraph{Answer to a).}
Take the matrix element of Eq.~(\ref{eq:7-eom}) with
$\bracket{0}{\delta\hat{Q}\,\cdot}{\,}$:
$\bracket{0}{\delta\hat{Q}[\hat{H},\hat{Q}^{\dagger}_{\nu}]}{0}
=\omega_{\nu}\bracket{0}{\delta\hat{Q}\hat{Q}^{\dagger}_{\nu}}{0}$.
Now use $\hat{Q}_{\nu}\bket{0}=0$, and its adjoint
$\bracket{0}{\hat{Q}^{\dagger}_{\nu}}{}=0$, to add for free the terms
$\bracket{0}{[\hat{H},\hat{Q}^{\dagger}_{\nu}]\delta\hat{Q}}{0}$ and
$\omega_{\nu}\bracket{0}{\hat{Q}^{\dagger}_{\nu}\delta\hat{Q}}{0}$, which
vanish because $\delta\hat{Q}$ is itself a de-excitation operator.
Subtracting turns both sides into commutators and gives
Eq.~(\ref{eq:7-eomdouble}).  The symmetrised form is preferable because the
double commutator can be evaluated with fewer contractions, and because it is
manifestly Hermitian.

\paragraph{Answer to b).}
With $\delta\hat{Q}=\hat{a}^{\dagger}_{i}\hat{a}_{m}$ and
$\hat{Q}^{\dagger}=\sum_{nj}X_{nj}\hat{a}^{\dagger}_{n}\hat{a}_{j}$ the
right-hand side of Eq.~(\ref{eq:7-eomdouble}) is
$\sum_{nj}X_{nj}\bracket{\mathrm{HF}}{[\hat{a}^{\dagger}_{i}\hat{a}_{m},
\hat{a}^{\dagger}_{n}\hat{a}_{j}]}{\mathrm{HF}}=X_{mi}$ by
Eq.~(\ref{eq:7-phcommutator}), exactly and without approximation, because the
expectation value of $\hat{a}^{\dagger}_{i}\hat{a}_{j}$ in a determinant is
$\delta_{ij}$ for occupied indices and zero otherwise.  The left-hand side is
$\sum_{nj}A_{mi,nj}X_{nj}$ with $A$ given by
Eq.~(\ref{eq:7-Amatrix}); evaluating the double commutator with the
generalised Wick theorem of Section~\ref{sec:wickgeneral} gives the one-body
piece $(\varepsilon_m-\varepsilon_i)\delta_{mn}\delta_{ij}$ from $\hat{H}_0$
and the single surviving contraction $\bar{v}_{mjin}$ from the interaction.

\paragraph{Answer to c).}
The $1p$-$1h$ determinants are orthonormal, and the Condon-Slater rules of
Section~\ref{sec:wickapps} give
\begin{equation}
  \bracket{\Phi^m_i}{\hat{H}}{\Phi^n_j}
  = E_{\mathrm{HF}}\delta_{mn}\delta_{ij}
  + \left(\varepsilon_m-\varepsilon_i\right)\delta_{mn}\delta_{ij}
  + \bar{v}_{mjin},
  \label{eq:7-exc1}
\end{equation}
which is $E_{\mathrm{HF}}\delta+A_{mi,nj}$.  Diagonalising $A$ is therefore
diagonalising $\hat{H}$ in the space spanned by $\bket{\mathrm{HF}}$ and all
singles -- CIS -- and reading the eigenvalues relative to the reference.  The
reference itself decouples by Brillouin's theorem, which is why it does not
appear in $A$.  The program checks
$\max|A-(\bm{H}_{1p1h}-E_{\mathrm{HF}})| = 3\times10^{-15}$.

\paragraph{Answer to d).}
Now $\delta\hat{Q}$ ranges over both
$\hat{a}^{\dagger}_{i}\hat{a}_{m}$ and $\hat{a}^{\dagger}_{m}\hat{a}_{i}$, so
Eq.~(\ref{eq:7-eomdouble}) becomes a pair of coupled equations for $X$ and
$Y$.  The left-hand sides give $A$ and $B$ as in
Eqs.~(\ref{eq:7-Amatrix}) and~(\ref{eq:7-Bmatrix}).  The right-hand sides
require
$\bracket{\mathrm{RPA}}{[\hat{a}^{\dagger}_{i}\hat{a}_{m},
\hat{a}^{\dagger}_{n}\hat{a}_{j}]}{\mathrm{RPA}}$ in the \emph{correlated}
ground state, which we do not have.  Replacing it by its value in the
Hartree-Fock determinant, Eq.~(\ref{eq:7-quasiboson}), is the quasiboson
approximation, and it is the only approximation made.  It produces the norm
matrix $\bm{\eta}=\mathrm{diag}(\bm{1},-\bm{1})$; the minus sign, from the
second row, is what makes Eq.~(\ref{eq:7-rpaequation}) non-Hermitian.

\paragraph{Answer to e).}
Take the complex conjugate of Eq.~(\ref{eq:7-rpaequation}) and interchange the
two blocks of rows:
$A^{*}X^{*}+B^{*}Y^{*}=\omega^{*}X^{*}$ and
$-BY^{*}-AX^{*}=\omega^{*}Y^{*}$, that is
\begin{equation}
  \begin{pmatrix}A & B\\ -B^{*} & -A^{*}\end{pmatrix}
  \begin{pmatrix}Y^{*}\\ X^{*}\end{pmatrix}
  = -\omega^{*}\begin{pmatrix}Y^{*}\\ X^{*}\end{pmatrix}.
  \label{eq:7-exe2}
\end{equation}
So the spectrum is symmetric about zero.  If $\omega$ is real this is the
partner $-\omega$; if it is complex the four roots
$\pm\omega,\pm\omega^{*}$ appear together, which is why imaginary
frequencies arrive in pairs.

\subsection*{Exercise 3: RPA, stability and Thouless}
\label{ex:7-thouless}
\addcontentsline{toc}{subsection}{Exercise 3: RPA, stability and Thouless}

\begin{enumerate}
  \item[a)] Show that $\hat{M}_{\mathrm{RPA}}=\bm{\eta}
  \hat{M}_{\mathrm{stab}}$ with $\bm{\eta}=\mathrm{diag}(\bm{1},-\bm{1})$,
  and deduce Eq.~(\ref{eq:7-stabilityequivalence}).
  \item[b)] Reproduce Table~\ref{tab:rpastability} for the Lipkin model and
  show that the collective root behaves as
  $\omega=\varepsilon\sqrt{1-\chi^2}$ near the transition.
  \item[c)] Verify Eq.~(\ref{eq:7-closedform}) for the pure pairing model by
  computing $A$ and $B$ analytically in the particle-hole channel.
  \item[d)] Show that $\hat{Q}_{\nu}\bket{\mathrm{RPA}}=0$ is solved by
  Eq.~(\ref{eq:7-rpagroundstate}), and compare its structure with Thouless'
  theorem and with the coupled-cluster ansatz.
\end{enumerate}

\paragraph{Answer to a).}
Multiplying $\hat{M}_{\mathrm{stab}}$ on the left by $\bm{\eta}$ flips the
sign of its second block row, turning $(B^{*},A^{*})$ into
$(-B^{*},-A^{*})$, which is $\hat{M}_{\mathrm{RPA}}$.  If
$\hat{M}_{\mathrm{stab}}>0$ it has a positive-definite Hermitian square root
$\bm{S}$, and
$\bm{S}^{-1}\hat{M}_{\mathrm{RPA}}\bm{S}
=\bm{S}^{-1}\bm{\eta}\bm{S}^{2}\bm{S}
=\bm{S}\bm{\eta}\bm{S}$, which is Hermitian.  Similar matrices have the same
spectrum, so the RPA eigenvalues are real.  If $\hat{M}_{\mathrm{stab}}$ has a
negative eigenvalue the square root is no longer Hermitian and the argument
collapses; the roots may then be imaginary, and for a single negative
direction they are.

\paragraph{Answer to b).}
Section~\ref{sec:lipkininstability} gives
$\Delta+A=\varepsilon\bm{I}$ and $\bm{B}=V(\bm{J}-\bm{I})$, whose largest
eigenvalue is $V(N-1)=\varepsilon\chi$.  Since $A$ and $B$ commute in this
model, the RPA roots are
$\omega=\sqrt{A^2-B^2}$ evaluated eigenvalue by eigenvalue, and the collective
one is
\begin{equation}
  \omega = \sqrt{\varepsilon^2-\varepsilon^2\chi^2}
   = \varepsilon\sqrt{1-\chi^2}.
  \label{eq:7-exb3}
\end{equation}
This vanishes at $\chi=1$ and is imaginary beyond, reproducing
Table~\ref{tab:rpastability}: at $\chi=0.9$ it gives
$\sqrt{1-0.81}=0.43589$, exactly the tabulated value.  Note also the square
root singularity: the mode softens with infinite slope as the transition is
approached, which is the mean-field signature of a second-order transition.

\paragraph{Answer to c).}
With pure pairing the Hartree-Fock orbitals are the original ones, the
occupied levels shifted down by $g/2$, so $\varepsilon_m-\varepsilon_i$ for
the lowest particle-hole pair is $\xi+g/2$.  The interaction contributes
$\bar{v}_{mjin}$ to $A$ and $\bar{v}_{mnij}$ to $B$; for the pairing force
only the terms in which $m,n$ and $i,j$ are spin partners of the same level
survive, giving a diagonal $A=\xi+g/2$ and $B=g/2$ in the collective channel.
Then $\omega=\sqrt{A^2-B^2}=\sqrt{(\xi+g/2)^2-(g/2)^2}=\sqrt{\xi^2+g\xi}$,
which is Eq.~(\ref{eq:7-closedform}).  Since $\xi^2+g\xi>0$ for all $g>-\xi$,
the particle-hole channel never becomes unstable.

\paragraph{Answer to d).}
Write $\bket{\mathrm{RPA}}=\mathcal{N}e^{\hat{S}}\bket{\mathrm{HF}}$ with
$\hat{S}=\tfrac12\sum Z_{mi,nj}\hat{a}^{\dagger}_{m}\hat{a}_{i}
\hat{a}^{\dagger}_{n}\hat{a}_{j}$.  Since the particle-hole operators commute
among themselves in the quasiboson approximation,
$[\hat{a}^{\dagger}_{i}\hat{a}_{m},\hat{S}]
=\sum_{nj}Z_{mi,nj}\hat{a}^{\dagger}_{n}\hat{a}_{j}$, and imposing
$\hat{Q}_{\nu}e^{\hat{S}}\bket{\mathrm{HF}}=0$ gives
$Y^{\nu}=\bm{Z}\,\bm{X}^{\nu}$ in matrix form, hence
$\bm{Z}=\bm{Y}^{*}(\bm{X}^{*})^{-1}$.  The structure is the same as Thouless'
theorem, Eq.~(\ref{eq:6-thouless}), with $2p$-$2h$ operators in place of the
$1p$-$1h$ ones -- so the RPA ground state stands to the doubles as a Slater
determinant does to the singles.  It is also the coupled-cluster doubles
ansatz $e^{\hat{T}_2}\bket{\Phi_0}$ of Eq.~(\ref{eq:5-ccsd}), with the
amplitudes fixed by the RPA equations rather than determined by their own
nonlinear equations.  RPA is thus an approximation to coupled-cluster doubles,
and the comparison between the two is a standard way of exposing what the
quasiboson approximation costs.

\subsection*{Exercise 4: BCS}
\label{ex:7-bcs}
\addcontentsline{toc}{subsection}{Exercise 4: BCS}

\begin{enumerate}
  \item[a)] Derive Eq.~(\ref{eq:7-bcsenergy}), being explicit about the
  $p=q$ and $p\ne q$ contributions.
  \item[b)] Show that only the $q=r$ part of $\hat{V}_{\mathrm{ph}}$ has a
  non-zero expectation value in the BCS state, and hence that the problem
  depends on $g$ and $f$ only through $g+2f$.
  \item[c)] Derive the critical condition~(\ref{eq:7-gc}) and solve it for
  $S_c$ and $g_c$.
  \item[d)] Verify that the Bogoliubov quasiparticles~(\ref{eq:7-bogoliubov})
  annihilate the BCS state and obey the fermion anticommutation relations.
  \item[e)] Compare the BCS ground-state energy with the exact one through
  the transition and comment on the sharpness of the mean-field transition.
\end{enumerate}

\paragraph{Answer to a).}
$\bket{\mathrm{BCS}}$ is a product over levels, so
$\langle\hat{N}\rangle=2\sum_p v_p^2$ and
$\langle\hat{H}_0\rangle=2\sum_p\varepsilon_p v_p^2$ immediately.  For the
interaction, $\hat{P}_q$ acting on the product picks out the $v_q$ term and
leaves level $q$ empty; $\hat{P}^{\dagger}_p$ in the bra does the same at $p$.
For $p=q$ the two operations are on the same level and the overlap is
$v_p^2$; for $p\ne q$ the surviving overlap has $u_p$ at $p$ and $u_q$ at $q$,
giving $u_pv_pu_qv_q$.  Assembling,
$\sum_{pq}\langle\hat{P}^{\dagger}_p\hat{P}_q\rangle
=\sum_p v_p^2+\sum_{p\ne q}u_pv_pu_qv_q$, and multiplying by $-S$ gives
Eq.~(\ref{eq:7-bcsenergy}).

\paragraph{Answer to b).}
$\hat{a}_{q-}\hat{a}_{r+}$ with $q\ne r$ leaves an unpaired particle at each
of $q$ and $r$; the resulting state has two broken pairs and is orthogonal to
anything $\hat{P}^{\dagger}_p$ can produce from
$\bket{\mathrm{BCS}}$, which contains only paired levels.  For $q=r$,
$\hat{a}_{q-}\hat{a}_{q+}=\hat{P}_q$ exactly, so the term reduces to
$-\tfrac{f}{2}\sum_{pq}\hat{P}^{\dagger}_p\hat{P}_q$ and its Hermitian
conjugate contributes an equal amount.  The total is
$-f\sum_{pq}\hat{P}^{\dagger}_p\hat{P}_q$, which added to
$-\tfrac{g}{2}\sum_{pq}\hat{P}^{\dagger}_p\hat{P}_q$ gives the single
strength $S=g/2+f$, that is $g\to g+2f$.  The program confirms that
$(g,f)=(1.4,0)$, $(1.2,0.1)$, $(0.8,0.3)$ and $(0,0.7)$ give identical
energies and gaps.

\paragraph{Answer to c).}
Linearise about $\Delta=0$.  There $v_p^2=1$ for the two lowest levels and
zero above, so the self-energy shifts them to
$\tilde{\varepsilon}_p=\varepsilon_p-S$ for $p\le2$ and leaves
$\tilde{\varepsilon}_p=\varepsilon_p$ above.  With $\varepsilon_p=0,1,2,3$
this gives $(-S,\,1-S,\,2,\,3)$, and the chemical potential sits midway
between the highest occupied and lowest empty level,
$\lambda=(3-S)/2$.  The linearised gap equation
$1=(S/2)\sum_p|\tilde{\varepsilon}_p-\lambda|^{-1}$ becomes
\begin{equation}
  1 = \frac{S}{2}\left[\frac{2}{(3+S)/2}\cdot 1
    + \frac{2}{(1+S)/2}\cdot 1\right]\cdot\frac12\cdot 2
   = 2S\left[\frac{1}{3+S}+\frac{1}{1+S}\right],
  \label{eq:7-exc4}
\end{equation}
which is Eq.~(\ref{eq:7-gc}) with $\xi=1$.  Solving numerically gives
$S_c=0.53519$ and $g_c=2S_c=1.07037$ at $f=0$, matching the bisection in the
program.

\paragraph{Answer to d).}
$\hat{\alpha}_{p+}\bket{\mathrm{BCS}}
=(u_p\hat{a}_{p+}-v_p\hat{a}^{\dagger}_{p-})
\prod_q(u_q+v_q\hat{P}^{\dagger}_q)\bket{0}$.  Only the factor at level $p$
matters; acting on $(u_p+v_p\hat{a}^{\dagger}_{p+}\hat{a}^{\dagger}_{p-})$
gives $u_pv_p\hat{a}^{\dagger}_{p-}-v_pu_p\hat{a}^{\dagger}_{p-}=0$.  For the
anticommutators,
\begin{equation}
  \left\{\hat{\alpha}_{p+},\hat{\alpha}^{\dagger}_{p+}\right\}
   = u_p^2\left\{\hat{a}_{p+},\hat{a}^{\dagger}_{p+}\right\}
   + v_p^2\left\{\hat{a}^{\dagger}_{p-},\hat{a}_{p-}\right\}
   = u_p^2+v_p^2 = 1,
  \label{eq:7-exd4}
\end{equation}
the cross terms
vanishing because they pair two creation or two annihilation operators.  The
normalisation $u_p^2+v_p^2=1$ is exactly what makes the transformation
canonical.  The program checks
$\max\Vert\hat{\alpha}_a\bket{\mathrm{BCS}}\Vert=3\times10^{-17}$.

\paragraph{Answer to e).}
From Table~\ref{tab:allresults}, $E_{\mathrm{BCS}}=E_{\mathrm{HF}}$ exactly
below $g_c=1.07$ and departs from it above.  The exact energy, by contrast,
runs smoothly through $g_c$ with no feature of any kind: at $g=1.0$, $1.2$ and
$1.4$ it is $0.6355$, $0.2645$ and $-0.1370$, a perfectly regular sequence.
The mean-field transition is therefore an artefact of restricting the trial
state to a product form.  What is physical is the crossover -- the gradual
build-up of pair correlations -- and what is spurious is its sharpness, and
with it the broken symmetry.  For a large system the crossover narrows and the
mean-field picture becomes accurate; for four particles it does not.

\subsection*{Exercise 5: QRPA and the spurious mode}
\label{ex:7-qrpa}
\addcontentsline{toc}{subsection}{Exercise 5: QRPA and the spurious mode}

\begin{enumerate}
  \item[a)] Show that if $[\hat{H},\hat{F}]=0$ while the reference is not an
  eigenstate of $\hat{F}$, then the amplitudes~(\ref{eq:7-spurious}) built
  from $\hat{F}$ give an exact zero mode of the RPA equations.
  \item[b)] Explain why cranking with $\hat{H}'=\hat{H}-\lambda\hat{N}$
  places the spurious root at zero, and why the root may come out as a small
  \emph{imaginary} pair rather than a small real one.
  \item[c)] Reproduce Table~\ref{tab:spurious} and confirm the unit overlap.
  What would you conclude if the overlap were small?
  \item[d)] Classify the QRPA phonons at $g=2$ by
  $\langle\Delta\hat{N}\rangle$, and explain why the shift is not quantised
  at $\pm2$.
  \item[e)] Explain why $E^{\mathrm{QRPA}}_{\mathrm{corr}}$ overshoots so
  badly in Table~\ref{tab:allresults}, and suggest what would fix it.
\end{enumerate}

\paragraph{Answer to a).}
Because $[\hat{H},\hat{F}]=0$, the operator $\hat{F}$ generates a family of
states $e^{\imath\theta\hat{F}}\bket{\mathrm{ref}}$ all with the same energy.
If the reference is not an eigenstate of $\hat{F}$ these are genuinely
different states, so the energy surface has a flat direction, and the second
derivative -- which is what $\hat{M}_{\mathrm{stab}}$ is -- has a zero
eigenvalue along it.  Algebraically, substituting the
amplitudes~(\ref{eq:7-spurious}) into the left-hand side of
Eq.~(\ref{eq:7-eomdouble}) gives
$\bracket{\mathrm{ref}}{[\delta\hat{Q},[\hat{H},\hat{F}]]}{\mathrm{ref}}=0$
identically, so $\omega=0$.

\paragraph{Answer to b).}
The BCS state extremises $\hat{H}-\lambda\hat{N}$, not $\hat{H}$, so it is
$\hat{H}'$ whose energy surface is flat at the BCS point and whose second
derivative therefore has the exact zero eigenvalue.  Building the QRPA from
$\hat{H}$ instead displaces the spurious root to a small non-zero value and
contaminates the physical ones.  The root may emerge imaginary because the RPA
metric $\bm{\eta}$ is singular in the spurious direction: there
$\sum(|X|^2-|Y|^2)=0$, the mode cannot be normalised, and rounding errors of
either sign are amplified into a small $\pm\imath\epsilon$ pair.  This is a
property of the numerics, not of the physics.

\paragraph{Answer to c).}
\texttt{demo\_qrpa()} prints the table.  The overlap between the near-zero
eigenvector and the amplitudes~(\ref{eq:7-spurious}) is $1.000000$ throughout,
so the root is unambiguously the number mode.  Were the overlap small, the
conclusion would be the opposite and far more interesting: a genuine soft mode,
signalling that the BCS solution is about to become unstable against some
other deformation, exactly as the Lipkin collective root does at $\chi=1$.
The two cases look identical in the spectrum alone, which is why the overlap
test matters.

\paragraph{Answer to d).}
Running \texttt{classify\_modes} gives $\langle\Delta\hat{N}\rangle=0$ to
machine precision for the roots at $2.6015$ and $2.9212$, and
$\pm0.357$ for the pair at $3.7879$.  The shift is not $\pm2$ because
$\bket{\mathrm{BCS}}$ is not an eigenstate of $\hat{N}$: it is a wave packet
in particle number, and $\langle\hat{N}\rangle$ in the phonon state measures
the \emph{mean} of a shifted packet, not a quantum number.  What is exact and
meaningful is the vanishing of the shift for the modes that stay within the
$N$-particle system, which is enforced by symmetry.  Number projection would
restore the quantisation.

\paragraph{Answer to e).}
Two effects compound.  First, Eq.~(\ref{eq:7-rpacorrelation}) sums over
\emph{all} roots, and the two-quasiparticle space contains pairing vibrations
that describe the neighbouring $N\pm2$ systems; their zero-point energy has no
business in the correlation energy of the $N$-particle system.  Second, the
quasiboson approximation is worse here than in the particle-hole channel,
because the BCS occupations are fractional and the commutator of two
two-quasiparticle operators is correspondingly further from a $c$-number.  The
cure for the first is particle-number projection before or after variation;
the cure for the second is a renormalised QRPA that keeps the leading
correction to Eq.~(\ref{eq:7-quasiboson}).  Both are standard, and both are
essential for a system with only four particles; in a heavy nucleus the
relative number fluctuation $\Delta N/N$ is small and plain QRPA is adequate.

\subsection*{Exercise 6: numerical explorations}
\label{ex:7-numerical}
\addcontentsline{toc}{subsection}{Exercise 6: numerical explorations}

\begin{enumerate}
  \item[a)] Using \texttt{rpa.py}, plot the exact, TDA and RPA lowest
  excitation energies against $g$ for $f=0.05g$ and for $f=0.5g$, and comment
  on where the approximations fail.
  \item[b)] Repeat Table~\ref{tab:allresults} for $L=6$ and $N=6$.  Does
  RPA improve or deteriorate as the system grows?
  \item[c)] Compute the RPA correlation energy from
  Eq.~(\ref{eq:7-rpacorrelation}) and, separately, as the difference between
  the exact and Hartree-Fock energies.  Plot the ratio against $g$.
  \item[d)] Find, by sweeping $f$ at fixed $g$, the coupling at which the
  particle-hole stability matrix of this model first acquires a negative
  eigenvalue -- or show that it never does within the range you can reach.
  \item[e)] Restrict the QRPA correlation energy to the
  $\langle\Delta\hat{N}\rangle=0$ modes and compare with
  Table~\ref{tab:allresults}.  How much of the overshoot does this remove?
\end{enumerate}

\paragraph{Answer.}
These are open-ended; a few pointers.  In b) the RPA correlation energy is an
extensive quantity while the quasiboson error per phonon is not, so RPA
generally improves in relative terms as the level space grows -- the opposite
of truncated configuration interaction, whose size-consistency error grows
(Section~\ref{sec:sizeconsistency}).  In c) the ratio starts near one at small
$g$ and drifts past it as RPA begins to overshoot; the crossing point is a
useful single number for characterising where the method turns over.  In d)
the answer for this model is that it never does within any reasonable range,
since Eq.~(\ref{eq:7-closedform}) is real for all $g>-\xi$ and the
particle-hole term only stiffens the matrix; the instability is in the pairing
channel, which requires the BCS reference to see.  In e) removing the
pairing-vibration roots takes out a large part of the overshoot but not all of
it, since the quasiboson approximation remains; the residual is a fair
estimate of what number projection alone cannot fix.
